\chapter{时间序列、协整与状态估计：从单位根到 OU}
\label{ch:lower-stat-arb-model}

\begin{itemize}
  % Generated from curriculum/manifest.json; do not edit by hand.
\item 直接先修：上册第 9 章《概率收敛、极限定理与集中不等式》、上册第 11 章《随机过程、Markov、鞅、Poisson 与 Brownian》、上册第 12 章《时间序列与状态空间基础》、上册第 14 章《浮点数、数值线代与病态问题》、上册第 15 章《Monte Carlo、bootstrap 与方差缩减》、上册第 16 章《研究有效性、样本外与交易摩擦》、上册第 17 章《可复现研究审计》。

  \item 学习产出：把单位根、协整、OU、状态空间、regime filtering 与变点检测串成可估计模型，并将预测接入仓位、成交、成本和失效监控。
  \item 研究边界：合成路径是数学 oracle；公开快照只验证数据协议，任何统计套利结论都必须经过候选选择、多重检验和 walk-forward 台账。
\end{itemize}

\section{先区分平稳、单位根与共同趋势}

AR(1) 模型 $X_t=\phi X_{t-1}+\varepsilon_t$ 在 $|\phi|<1$ 且创新满足适当矩条件时具有稳定二阶结构。若 $\phi=1$，冲击永久进入水平，通常要研究 $\Delta X_t$。Dickey--Fuller 回归
\[
\Delta X_t=a+\gamma X_{t-1}+u_t
\]
检验 $\gamma=0$，其 t 统计量不服从普通 Student t 分布，必须使用与确定性项、滞后阶相匹配的临界值。ADF 不是“平稳证明器”：低功效、结构突变和错误滞后都可能改变结论。
所有滞后、变换与预测都必须可由决策时信息集 $\mathcal F_t$ 构造；全样本选规格不属于该信息集。

\begin{diagnostic}
先画水平、差分和滚动尺度，再确定确定性项与最大滞后。看到一个小 p 值以后才选择含趋势或不含趋势的规格，会把规格搜索藏进检验。
\end{diagnostic}

\MFQTransition{单位根检验的原假设与检验力}

ADF 回归在差分方程中检验水平项系数，原假设是单位根；KPSS 的原假设则是平稳。二者结论可以组合成“平稳”“单位根”或“证据不足”。短样本、接近单位根和结构突变会显著降低检验力，不能把 $p>0.05$ 翻译成“证明存在单位根”。

确定性项必须事前决定。对带趋势序列漏掉趋势会错误拒绝或接受；滞后阶数过低残差相关，过高又损失自由度。价格通常先看对数水平与差分，利差和估值比率则要依据经济含义选择变换。

\section{Engle--Granger 与误差修正}

若 $x_t,y_t$ 都是 $I(1)$，但存在 $\beta$ 使 $z_t=y_t-a-\beta x_t$ 为 $I(0)$，则二者协整。Engle--Granger 两步法先估长期关系，再对残差做专用单位根检验\cite{engle1987cointegration}。用于核对运算的确定性算例令
\[
y_t=1.5+2x_t+0.55^t,
\]
对此数据作 OLS 得斜率约 $1.9912$，残差 DF 统计量约 $-8.3135$。斜率没有精确等于 2，是有限样本中确定性残差与趋势相关造成的差异。这里 $0.55^t$ 没有随机创新，不能用这一条序列说明平稳模型的检验力。

配套教学实验另使用 $x_t=x_{t-1}+u_t$、$z_t=\phi z_{t-1}+\eta_t$、$y_t=a+\beta x_t+z_t$，其中 $u_t\sim N(0,1)$、$\eta_t\sim N(0,\sigma_\eta^2)$ 相互独立且跨期独立，$|\phi|<1$。令 $z_0\sim N(0,\sigma_\eta^2/(1-\phi^2))$ 并独立于未来创新，则价差从一开始就平稳，而两个水平序列共享随机趋势。改变样本量并重复估计，才能看到有限样本误差。

% Generated by tools/build_figures.py; edit figures/figure-specs.json instead.
\MFQFigure{tex/figures/generated/lower-sa-cointegration.pdf}{两条非平稳价格路径可以共享共同趋势，而特定线性组合保持平稳并围绕均值波动。}{lower-sa-cointegration}





长期偏离进入误差修正模型：
\[
\Delta y_t=c+\alpha z_{t-1}+\delta\Delta x_t+u_t.
\]
合成例得到 $\hat\alpha\approx-0.4606$；负号表示 $y_t$ 对正偏离向下调整。只有在协整向量、方向和时点协议固定后，$\alpha$ 才能解释为调整速度。

\section{选读：Johansen 的多变量秩问题}

VAR 的 VECM 形式为
\[
\Delta w_t=\Pi w_{t-1}+\sum_{j=1}^{p-1}\Gamma_j\Delta w_{t-j}+u_t,
\qquad \Pi=\alpha\beta^{\mathsf T}.
\]
$\operatorname{rank}(\Pi)=r$ 是协整秩。Johansen 方法用广义特征值构造 trace 与最大特征根统计量，并按确定性项与滞后阶读取非标准临界值\cite{johansen1991estimation}。本项目的透明奇异值诊断只帮助理解“共同趋势方向”；正式 rank=1 结论由 \texttt{statsmodels} Johansen 临界值表独立计算。两者假设不同，不能把透明诊断冒充正式检验。

\section{从离散 AR 映射到 OU}

连续时间 OU 过程
\[
dZ_t=\kappa(\mu-Z_t)dt+\sigma dW_t
\]
对 $e^{\kappa t}(Z_t-\mu)$ 使用乘积法则，再在 $[t,t+\Delta]$ 积分，得到
\[
Z_{t+\Delta}=\mu+e^{-\kappa\Delta}(Z_t-\mu)
+\sigma\int_t^{t+\Delta}e^{-\kappa(t+\Delta-s)}dW_s.
\]
最后一项是独立于时点 $t$ 信息的零均值正态创新，方差为 $\sigma^2(1-e^{-2\kappa\Delta})/(2\kappa)$，所以 $\phi=e^{-\kappa\Delta}$。条件均值偏离按指数衰减，令其减半即得到半衰期。当 $0<\hat\phi<1$ 时，
\[
\hat\kappa=-\frac{\log\hat\phi}{\Delta},
\qquad t_{1/2}=\frac{\log 2}{\hat\kappa}.
\]
合成残差给出 $\hat\phi\approx0.5442$、半衰期约 $1.1393$ 个周期。$1/\hat\kappa$ 只是特征回归时间，不是首次通过时间。若从 $z_0>0$ 首次命中均值边界 0，扩散率为 $\sigma$，边值问题的解为
\[
\mathbb E_{z_0}\tau_0
=\frac{\sqrt\pi}{\kappa}
\int_0^{z_0\sqrt\kappa/\sigma}e^{u^2}\operatorname{erfc}(u)\,du.
\]
实现从离散创新方差映射 $\sigma$，并数值积分该式；起点、边界或扩散率改变时必须重算，不能用半衰期替代。

% Generated by tools/build_figures.py; edit figures/figure-specs.json instead.
\MFQFigure{tex/figures/generated/lower-sa-ou.pdf}{OU 漂移强度决定条件均值回归速度；半衰期描述偏离缩小一半所需的时间尺度。}{lower-sa-ou}





\begin{note}[Python 边界]
\texttt{adfuller}、\texttt{coint} 与 \texttt{coint\_johansen} 的默认确定性项、自动滞后和临界值并不相同。调用函数前先把统计问题写成完整规格，再在报告中保存库版本与参数。
\end{note}

\paragraph{实验阅读。}先比较确定性衰减与随机平稳价差，再计算 OLS 残差、协整检验和误差修正系数。公开宏观快照用于理解观测频率与真实数据背景，不作为可交易配对。

\section{ARMA、VAR 与 GARCH 的职责}

ARMA 描述条件均值的线性动态，VAR 处理多变量相互滞后，GARCH 描述条件方差：
\[
r_t=\mu_t+\varepsilon_t,\qquad
\varepsilon_t=\sigma_tz_t,\qquad
\sigma_t^2=\omega+\alpha\varepsilon_{t-1}^2+\beta\sigma_{t-1}^2.
\]
$\alpha+\beta$ 接近 1 表示波动冲击衰减慢；它不是价格均值回归速度。把 GARCH 波动预测直接当方向信号会混淆条件均值与条件方差。

残差诊断至少包括自相关、平方残差自相关、尾部分布和参数稳定。模型阶数可用训练期信息准则筛选，再用滚动样本外验证；不能在全样本 ACF 图上挑阶数后仍称测试独立。

\MFQTransition{协整向量的经济解释与不稳定性}

两个 $I(1)$ 序列若存在 $\beta$ 使 $z_t=y_t-\beta x_t$ 平稳，则共享随机趋势。$\beta$ 的尺度要固定，例如把 $y$ 系数规范为 1。纯统计搜索会从大量资产中找到偶然平稳组合，因此候选形成机制和多重检验必须记录。

滚动 $\hat\beta_t$ 变化意味着对冲比率变化，也可能说明关系失效。若每期重新估计并立刻用新 beta 重写过去价差，历史信号会被回看修订；研究账本要保存当时版本的 beta 和价差。

这组模型为后续路线 Capstone 提供状态、半衰期与失效边界；它们描述的是可检验的动态假设，不是收益承诺。
